\chapter{Concrete Examples}\label{ch:examples}

This chapter collects concrete MPS tensors from
\cite[Section~III]{Cirac2021Matrix}.  Each example is specified by its
physical dimension~$d$, bond dimension~$D$, and explicit matrix
family~$\{A^i\}$.  Key structural properties (injectivity, RFP status,
symmetry) are stated for each tensor.

%% ===================================================================
\section{The GHZ state}\label{sec:ghz}
%% ===================================================================

The Greenberger--Horne--Zeilinger (GHZ) state is the prototypical
non-injective, renormalization-fixed-point MPS.

\begin{definition}[GHZ tensor]\label{def:ghz_tensor}
    \lean{MPSTensor.ghzTensor}
    \leanok
    \uses{def:mps_tensor}
    Set $d = D = 2$.  The GHZ tensor is the family
    $\{A^0, A^1\} \subset \MN{2}$ defined by
    \[
        A^0 = \ketbra{0}{0} = \diag(1,0),
        \qquad
        A^1 = \ketbra{1}{1} = \diag(0,1).
    \]
    For a system of $N$ sites the resulting MPV is the GHZ state
    \[
        \ket{\mathrm{GHZ}_N}
        = \ket{0}^{\otimes N} + \ket{1}^{\otimes N}.
    \]
\end{definition}

\begin{theorem}[GHZ transfer map is idempotent]\label{thm:ghz_transfer_idempotent}
    \lean{MPSTensor.ghz_transferMap_idempotent}
    \leanok
    \uses{def:ghz_tensor, def:transfer_map}
    The transfer map of the GHZ tensor satisfies $\E_A^2 = \E_A$.
\end{theorem}

\begin{proof}
    \leanok
    The transfer map acts by
    $\E_A(X)_{ij} = \delta_{ij} X_{ij}$
    (extraction of the diagonal).
    This operation is idempotent.
\end{proof}

\begin{theorem}[GHZ tensor is an RFP]\label{thm:ghz_is_rfp}
    \lean{MPSTensor.ghz_isRFP}
    \leanok
    \uses{def:ghz_tensor, def:is_rfp, thm:ghz_transfer_idempotent}
    The GHZ tensor is a renormalization fixed point.
\end{theorem}

\begin{proof}
    \leanok
    \uses{thm:ghz_transfer_idempotent}
    Immediate from Theorem~\ref{thm:ghz_transfer_idempotent}.
\end{proof}

\begin{theorem}[GHZ tensor is not injective]\label{thm:ghz_not_injective}
    \lean{MPSTensor.ghz_not_isInjective}
    \leanok
    \uses{def:ghz_tensor, def:injective}
    The GHZ tensor is not injective: $\spn_{\C}\{A^0, A^1\}$ is the
    two-dimensional space of diagonal matrices in~$\MN{2}$, which does
    not equal~$\MN{2}$.
\end{theorem}

\begin{proof}
    \leanok
    Every element of $\spn_{\C}\{A^0, A^1\}$ has zero off-diagonal
    entries, so the matrix with all entries equal to~$1$ is not in the
    span.
\end{proof}

\begin{theorem}[$\Z_2$ on-site symmetry of the GHZ tensor]\label{thm:ghz_z2_symmetric}
    \lean{MPSTensor.ghz_isOnSiteSymmetric_Z2}
    \leanok
    \uses{def:ghz_tensor, def:on_site_symmetric}
    Let $\Z_2 = \{1, \sigma_x\}$ act on $\C^2$ by the Pauli-$X$
    matrix~$\sigma_x$.  The GHZ tensor is on-site symmetric under this
    action: for each $g \in \Z_2$, the twisted tensor $\widetilde A_g$
    defines the same MPV family as $A$, i.e.\ $\mc V(A)=\mc V(\widetilde
        A_g)$.
\end{theorem}

\begin{proof}
    \leanok
    The identity element is trivial.  For the generator
    $g = \sigma_x$, the twisted tensor satisfies
    $\widetilde A_g^i = A^{1-i}$, i.e.\ the two matrices are swapped.
    Conjugation by $\sigma_x$ implements the same swap, giving
    gauge equivalence and hence the same MPV family.
\end{proof}

%% ===================================================================
\section{The AKLT state}\label{sec:aklt}
%% ===================================================================

The Affleck--Kennedy--Lieb--Tasaki (AKLT) state is the canonical
example of a normal (block-injective) MPS with non-trivial
symmetry-protected topological (SPT) order.

\begin{definition}[AKLT tensor]\label{def:aklt_tensor}
    \lean{MPSTensor.akltTensor}
    \leanok
    \uses{def:mps_tensor}
    Set $d = 3$ (spin-1) and $D = 2$.  The AKLT tensor is the family
    $\{A^0, A^1, A^2\} \subset \MN{2}$ defined via the Pauli matrices by
    \[
        A^0 = \frac{1}{\sqrt{3}}\,\sigma_z,
        \qquad
        A^1 = \frac{\sqrt{2}}{\sqrt{3}}\,\ketbra{0}{1},
        \qquad
        A^2 = -\frac{\sqrt{2}}{\sqrt{3}}\,\ketbra{1}{0},
    \]
    where
    $\sigma_z = \ketbra{0}{0} - \ketbra{1}{1}$,
    so equivalently $A^1 = \frac{1}{\sqrt{3}}\,\sigma_+$ and
    $A^2 = -\frac{1}{\sqrt{3}}\,\sigma_-$ for
    $\sigma_+ = \sqrt{2}\,\ketbra{0}{1}$ and
    $\sigma_- = \sqrt{2}\,\ketbra{1}{0}$.
    Equivalently, the matrices are the Clebsch--Gordan coefficients
    coupling spin-$1$ and spin-$\tfrac12$ to spin-$\tfrac12$.
\end{definition}

\begin{theorem}[AKLT tensor is not injective]\label{thm:aklt_not_injective}
    \lean{MPSTensor.aklt_not_isInjective}
    \leanok
    \uses{def:aklt_tensor, def:injective}
    The AKLT tensor is not injective at a single site:
    $\spn_{\C}\{A^0, A^1, A^2\}$ is the three-dimensional space of
    traceless matrices $\mathfrak{sl}(2,\C)$, which is a proper
    subspace of~$\MN{2}$.
\end{theorem}

\begin{theorem}[AKLT tensor is normal]\label{thm:aklt_normal}
    \lean{MPSTensor.aklt_isNormal}
    \leanok
    \uses{def:aklt_tensor, def:transfer_map}
    The transfer map of the AKLT tensor has a unique eigenvalue of
    maximal modulus (a simple spectral-radius eigenvalue); in
    particular the AKLT tensor is normal.  Concretely, the
    length-$2$ blocked tensor is injective.
\end{theorem}

\begin{theorem}[$\Z_2 \times \Z_2$ on-site symmetry of the AKLT tensor]\label{thm:aklt_z2z2_symmetric}
    \lean{MPSTensor.aklt_isOnSiteSymmetric_Z2Z2}
    \leanok
    \uses{def:aklt_tensor, def:on_site_symmetric}
    The AKLT tensor is on-site symmetric under the $\Z_2 \times \Z_2$ subgroup
    of $\mathrm{SO}(3)$ generated by two commuting $\pi$-rotations about
    orthogonal axes, acting on the spin-$1$ physical index.
\end{theorem}

\begin{theorem}[Anticommuting virtual gauges of the AKLT symmetry]\label{thm:aklt_gauge_anticomm}
    \lean{MPSTensor.aklt_gauge_anticomm}
    \leanok
    The two virtual matrices that implement the symmetry of
    Theorem~\ref{thm:aklt_z2z2_symmetric}, namely $i\sigma_y$ and $\sigma_z$,
    anticommute.
\end{theorem}

\begin{remark}\label{rem:aklt_z2z2_spt}
    Because the physical generators commute while their virtual representatives
    anticommute (Theorem~\ref{thm:aklt_gauge_anticomm}), the bond representation
    is projective, with $2$-cocycle the non-trivial element of
    $H^2(\Z_2 \times \Z_2, \mathrm{U}(1)) \cong \Z_2$, the signature of a
    non-trivial SPT phase.  This $\Z_2 \times \Z_2$ is the smallest subgroup of
    $\mathrm{SO}(3)$ that still carries a non-trivial $2$-cocycle.
\end{remark}

\begin{definition}[Cartesian AKLT tensor]\label{def:aklt_cartesian}
    \lean{MPSTensor.akltCartesian}
    \leanok
    \uses{def:mps_tensor}
    The Cartesian presentation of the AKLT tensor has $d = 3$, $D = 2$, and
    $A^x = \sigma_x$, $A^y = \sigma_y$, $A^z = \sigma_z$, the three Pauli
    matrices.  This is the rotationally natural form used in
    \cite{Cirac2021Matrix} (around line~1159); it is the AKLT state of
    Definition~\ref{def:aklt_tensor} written in the Cartesian basis of the
    spin-$1$ site.
\end{definition}

\begin{theorem}[The spin-$\tfrac12$ cover is surjective onto $\mathrm{SO}(3)$]\label{thm:su2_so3_surjective}
    \lean{MPSTensor.spinHalfCover_surjective_onto_SO3}
    \leanok
    \uses{def:aklt_cartesian}
    Every rotation $R \in \mathrm{SO}(3)$ is realized by conjugating the Pauli
    vector by some $U \in \mathrm{SU}(2)$: there is $U \in \mathrm{SU}(2)$ with
    $\bigl(\tfrac12 \tr(\sigma_i\, U \sigma_j U^{-1})\bigr)_{ij} = R$.
    Equivalently, the adjoint double cover $\mathrm{SU}(2)\to\mathrm{SO}(3)$ is
    surjective.  The proof factors $R$ by its $ZXZ$ Euler decomposition and hits
    each coordinate rotation with an explicit $\mathrm{SU}(2)$ element.
\end{theorem}

\begin{theorem}[$\mathrm{SO}(3)$ on-site symmetry of the AKLT tensor]\label{thm:aklt_so3_symmetric}
    \lean{MPSTensor.aklt_isOnSiteSymmetric_SO3}
    \leanok
    \uses{def:aklt_cartesian, def:on_site_symmetric, thm:su2_so3_surjective}
    The Cartesian AKLT tensor is on-site symmetric under the spin-$1$ (defining)
    representation of $\mathrm{SO}(3)$ on the physical index.
\end{theorem}

\begin{remark}\label{rem:aklt_so3_spt}
    The bond index carries the projective spin-$\tfrac12$ representation: the
    gauge for a rotation is either of its two $\mathrm{SU}(2)$ preimages, so the
    assignment is a representation only up to sign and factors through the double
    cover $\mathrm{SU}(2)\to\mathrm{SO}(3)$.  Its class is the non-trivial element
    of $H^2(\mathrm{SO}(3),\mathrm{U}(1))\cong\Z_2$, witnessing the SPT phase.
    This cohomological classification, which needs the continuous group
    cohomology of $\mathrm{SO}(3)$, is not formalized; the obstruction is
    recorded in \texttt{docs/paper-gaps/rmp\_aklt\_continuous\_rotation\_gap.tex}.
\end{remark}

%% ===================================================================
\section{The cluster state}\label{sec:cluster}
%% ===================================================================

The cluster state is a universal resource for measurement-based quantum
computation and provides another example of a non-trivial SPT phase.

\begin{definition}[Cluster-state tensor]\label{def:cluster_tensor}
    \lean{MPSTensor.clusterTensor}
    \leanok
    \uses{def:mps_tensor}
    Set $d = D = 2$.  The cluster-state tensor is the family
    $\{A^0, A^1\} \subset \MN{2}$ defined by
    \[
        A^0 = \ketbra{+}{0}
        = \frac{1}{\sqrt{2}}
        \begin{pmatrix} 1 & 0 \\ 1 & 0 \end{pmatrix},
        \qquad
        A^1 = \ketbra{-}{1}
        = \frac{1}{\sqrt{2}}
        \begin{pmatrix} 0 & 1 \\ 0 & -1 \end{pmatrix},
    \]
    where $\ket{\pm} = \tfrac{1}{\sqrt{2}}(\ket{0} \pm \ket{1})$.
    The review~\cite{Cirac2021Matrix} writes the reflected convention
    $A^0 = \ketbra{0}{+}$, $A^1 = \ketbra{1}{-}$; the transpose taken
    here is the same state read in the opposite direction and carries the
    same SPT order.
\end{definition}

\begin{theorem}[Cluster-state tensor is not injective]\label{thm:cluster_not_injective}
    \lean{MPSTensor.cluster_not_isInjective}
    \leanok
    \uses{def:cluster_tensor, def:injective}
    The cluster-state tensor is not injective:
    $\spn_{\C}\{A^0, A^1\} = \spn_{\C}\{\ketbra{+}{0}, \ketbra{-}{1}\}$
    has dimension~$2$, not $\dim \MN{2} = 4$.
\end{theorem}

\begin{proof}
    \leanok
    \uses{def:cluster_tensor}
    Both $A^0$ and $A^1$ are rank-one matrices with linearly
    independent column spaces, so they are linearly independent.
    Their span is two-dimensional, which is strictly less than
    $\dim \MN{2} = 4$.
\end{proof}

\begin{theorem}[Cluster-state tensor is normal]\label{thm:cluster_normal}
    \lean{MPSTensor.cluster_isNBlkInjective_two}
    \leanok
    \uses{def:cluster_tensor, def:injective}
    Although the cluster-state tensor is not injective at a single site, its
    length-$2$ blocking is injective: the four length-$2$ products span
    $\MN{2}$.  In particular the cluster-state tensor is normal.
\end{theorem}

\begin{remark}\label{rem:cluster_injective_blocked}
    Although the cluster-state tensor is not injective at length~$1$,
    the length-$2$ blocked tensor is injective.  Using
    $\braket{0}{\pm} = \tfrac{1}{\sqrt{2}}$,
    $\braket{1}{+} = \tfrac{1}{\sqrt{2}}$, and
    $\braket{1}{-} = -\tfrac{1}{\sqrt{2}}$, the
    four products $A^{ij} = A^i A^j$ are
    \[
        A^{00} = \tfrac{1}{\sqrt{2}}\ketbra{+}{0},\quad
        A^{01} = \tfrac{1}{\sqrt{2}}\ketbra{+}{1},\quad
        A^{10} = \tfrac{1}{\sqrt{2}}\ketbra{-}{0},\quad
        A^{11} = -\tfrac{1}{\sqrt{2}}\ketbra{-}{1}.
    \]
    Since $\{\ket{+},\ket{-}\}$ and $\{\bra{0},\bra{1}\}$ are bases,
    these four rank-one matrices are linearly independent and
    span~$\MN{2}$.
\end{remark}

\begin{theorem}[$\Z_2 \times \Z_2$ on-site symmetry of the blocked cluster-state tensor]\label{thm:cluster_z2z2_symmetric}
    \lean{MPSTensor.cluster_isOnSiteSymmetric_Z2Z2}
    \leanok
    \uses{def:cluster_tensor, def:on_site_symmetric}
    The length-$2$ blocked cluster-state tensor is on-site symmetric
    under $\Z_2 \times \Z_2$.  The two generators act on the
    $4$-dimensional blocked physical space $(\C^2)^{\otimes 2}$ by
    $\sigma_x \otimes I$ and $I \otimes \sigma_x$,
    which commute and each square to the identity, defining a linear
    representation of $\Z_2 \times \Z_2$.
\end{theorem}

\begin{theorem}[Anticommuting virtual gauges of the cluster-state symmetry]\label{thm:cluster_gauge_anticomm}
    \lean{MPSTensor.cluster_gauge_anticomm}
    \leanok
    The virtual matrices implementing the two generators of
    Theorem~\ref{thm:cluster_z2z2_symmetric}, namely $V_1 = \sigma_z$ and
    $V_2 = \sigma_x$, anticommute ($V_1 V_2 = -V_2 V_1$).
\end{theorem}

\begin{remark}\label{rem:cluster_z2z2_spt}
    Since the physical generators commute while their virtual representatives
    anticommute (Theorem~\ref{thm:cluster_gauge_anticomm}), the bond
    representation is projective, with $2$-cocycle the non-trivial element of
    $H^2(\Z_2 \times \Z_2, \mathrm{U}(1)) \cong \Z_2$, witnessing a non-trivial
    SPT phase.
\end{remark}
